\section{Conclusion}
\label{sec:conclusion}

Geographic resolution is a dial, not a binary choice.
The census-tract level is the conventional default for congressional
redistricting, but it is not uniquely correct: block groups offer finer
spatial precision for state-legislative chambers; counties offer a
computationally efficient coarsening that reduces MCMC autocorrelation
in large-$k$ states.
Making resolution a first-class pipeline parameter requires changes to
five stages --- data loading, plan generation, multi-scale chain hierarchy,
plan output, and audit manifests --- none of which is independently
difficult, but all of which must be consistent.

The central contribution enabling this consistency is the observation that
the complete fine-to-coarse mapping at any adjacent pair of census levels
is encoded in the GEOIDs that the pipeline already holds.
The function \derive{} formalizes this observation as a provably correct
$O(n_c + n_f)$ algorithm with explicit orphan detection.
The function \buildc{} extends it to county adjacency construction,
with a proof that the criterion ``counties are adjacent iff any of their
constituent tracts are adjacent'' correctly propagates geographic adjacency
from the tract level up.

Option B (tract$\to$county multi-scale) is fully deployed, requires no
additional data, and reduces lag-100 Hamming autocorrelation by
approximately 27\%\est\ on Texas $k=38$ relative to single-scale ReCom
at tract resolution\footnote{\est Estimated value from a single 2000-step run,
seed $s=42$. Phase~2 will provide multi-run validation with \texttt{criterion.rs}
benchmarks.}.
Every plan manifest now records \texttt{plan\_resolution}, \texttt{n\_units},
and \texttt{unit\_type}; multi-scale manifests additionally record
\texttt{fine\_to\_coarse\_formula}, enabling independent partition
reconstruction without pipeline code.

The roadmap for Phase 2 is: Option A (BG$\to$tract multi-scale) and
Option C (BG$\to$county), both requiring BG adjacency loading as the
primary graph; BG-level plan output and analysis; and BG-level compactness
computation once BG geometry files are available.

Resolution is not an implementation detail.
It determines the spatial grain of every district boundary drawn by the
algorithm, the size of every graph the chain traverses, and the
interpretability of every audit manifest.
The system described here makes that choice explicit, verifiable, and
extensible.

\paragraph{Acknowledgments.}
Census data from the U.S. Census Bureau, P.L.\ 94-171 Redistricting Data.
All computations performed with the \redist{} pipeline~\citep{dellaLibera2026b1}.

\paragraph{Code availability.}
The implementation is in \texttt{redist-cli/src/adjacency\_loader.rs}
(\derive{}, \buildc{}) and \texttt{redist-core/src/manifest.rs}
(\texttt{plan\_resolution} fields).
The full pipeline is available at the project repository.
